\documentclass[11pt]{article}
\usepackage{amsmath}
\usepackage{amssymb}
\usepackage{amsthm}
\usepackage[margin=1in]{geometry}

\newtheorem{theorem}{Theorem}
\newtheorem{lemma}[theorem]{Lemma}
\newtheorem{corollary}[theorem]{Corollary}
\theoremstyle{definition}
\newtheorem{definition}[theorem]{Definition}
\newtheorem{remark}[theorem]{Remark}

\title{Refutation of Conjecture 00000001367:\\[2pt]
The Borel Chromatic Number of the $\mathbb{Z}^2$ Grid is $2$, Not $5$}
\author{}
\date{\today}

\begin{document}
\maketitle

\begin{abstract}
We refute conjecture 00000001367, which asserts that the Borel chromatic number
of the Cayley graph of $\mathbb{Z}^2$ with generators $\{\pm e_1,\pm e_2\}$ is
exactly $5$ while its classical chromatic number is $4$, the Borel constraint
allegedly forcing one extra color. The graph in question is the square grid,
which is bipartite; hence its classical chromatic number is $2$, not $4$. Since
$\mathbb{Z}^2$ is countable and carries the discrete $\sigma$-algebra, every
subset of the vertex set is Borel, so the parity $2$-coloring is Borel; hence
the Borel chromatic number is also $2$, not $5$. Both asserted values are wrong,
and the Borel constraint forces no extra color here.
\end{abstract}

\section{The conjecture}

The conjecture under consideration reads verbatim:

\begin{quote}
\emph{Definition: The Borel chromatic number of a Borel graph on a standard Borel
space is the least number of colors in a Borel-measurable proper coloring.
Conjecture: The Borel chromatic number of the Cayley graph of $\mathbb{Z}^2$
(generators $\pm e_1,\pm e_2$) is exactly $5$; the classical chromatic number is
$4$, but the Borel constraint forces one extra color.}
\end{quote}

We show that both numerical assertions ($4$ and $5$) are false.

\section{The graph is the square grid}

\begin{definition}[Cayley graph]
Let $G$ be a group and $S\subseteq G$ a symmetric generating set (i.e.\
$S=S^{-1}$ and $e\notin S$). The \emph{Cayley graph} $\operatorname{Cay}(G,S)$
has vertex set $G$ and an (undirected) edge between $g$ and $h$ iff
$h=gs$ for some $s\in S$.
\end{definition}

Take $G=\mathbb{Z}^2$ (written additively) and $S=\{\pm e_1,\pm e_2\}$, where
$e_1=(1,0)$ and $e_2=(0,1)$. Since $G$ is abelian and $S=-S$, the edge relation
becomes: $p,q\in\mathbb{Z}^2$ are adjacent iff
\[
q=p+e_1,\quad q=p-e_1,\quad q=p+e_2,\quad\text{or}\quad q=p-e_2.
\]
Equivalently, writing $p=(i,j)$ and $q=(i',j')$, two lattice points are adjacent
iff $|i-i'|+|j-j'|=1$. This is precisely the \emph{square grid}.

\section{Bipartiteness and the classical chromatic number}

\begin{definition}[Proper coloring, chromatic number]
A \emph{proper coloring} of a graph $G=(V,E)$ is a map $c\colon V\to C$ such
that $c(p)\neq c(q)$ whenever $\{p,q\}\in E$. The \emph{chromatic number}
$\chi(G)$ is the least cardinality $|C|$ for which a proper coloring exists.
\end{definition}

\begin{theorem}[Bipartiteness]
\label{thm:bipartite}
The square grid $G=\operatorname{Cay}(\mathbb{Z}^2,\{\pm e_1,\pm e_2\})$ is
bipartite. Explicitly, the map
\[
c_0\colon \mathbb{Z}^2\to\{0,1\},\qquad c_0(i,j) := (i+j)\bmod 2,
\]
is a proper $2$-coloring.
\end{theorem}

\begin{proof}
Let $p=(i,j)$ and let $q$ be adjacent to $p$. There are four cases.
\begin{itemize}
\item $q=(i+1,j)$: then $q_1+q_2 = i+j+1$, so
$c_0(q)\equiv (i+j+1)\equiv c_0(p)+1 \pmod 2$.
\item $q=(i-1,j)$: then $q_1+q_2 = i+j-1\equiv c_0(p)+1 \pmod 2$.
\item $q=(i,j+1)$: then $q_1+q_2=i+j+1\equiv c_0(p)+1\pmod 2$.
\item $q=(i,j-1)$: then $q_1+q_2=i+j-1\equiv c_0(p)+1\pmod 2$.
\end{itemize}
In every case $c_0(q)\equiv c_0(p)+1 \pmod 2$, hence $c_0(q)\neq c_0(p)$. Thus
$c_0$ is proper, and the bipartition is given by the two parity classes of
$i+j$.
\end{proof}

\begin{corollary}[Classical chromatic number]
\label{cor:chi}
$\chi(G)=2$. In particular $\chi(G)\neq 4$.
\end{corollary}

\begin{proof}
By Theorem~\ref{thm:bipartite}, $c_0$ is a proper $2$-coloring, so
$\chi(G)\le 2$. On the other hand $G$ has at least one edge, for instance
$(0,0)$--$(1,0)$; no graph with an edge admits a proper $1$-coloring, so
$\chi(G)\ge 2$. Hence $\chi(G)=2$.
\end{proof}

\section{Countability and the Borel chromatic number}

We recall the relevant definitions.

\begin{definition}[Borel graph and Borel chromatic number]
A \emph{Borel graph} on a standard Borel space $(X,\mathcal{B})$ is a pair
$(X,E)$ where $E\subseteq X\times X$ is a symmetric, irreflexive Borel set. Its
\emph{Borel chromatic number} $\chi_B(G)$ is the least cardinality of a standard
Borel space $C$ admitting a Borel-measurable proper coloring $c\colon X\to C$.
\end{definition}

The vertex set $\mathbb{Z}^2$ is countable. We equip it with the discrete
$\sigma$-algebra $\mathcal{B}=\mathcal{P}(\mathbb{Z}^2)$, i.e.\ every subset is
measurable. The pair $(\mathbb{Z}^2,\mathcal{P}(\mathbb{Z}^2))$ is a standard
Borel space: it is Borel-isomorphic to a countable discrete Polish space (for
example, transport a bijection $\mathbb{Z}^2\to\mathbb{N}$). Because the space is
discrete, the edge set is all of $E\subseteq V\times V$, a countable union of
singletons, hence Borel, and $G$ is a Borel graph. The following lemma is the
key point.

\begin{lemma}[Every subset is Borel]
\label{lem:all-borel}
On the countable discrete space $\mathbb{Z}^2$, every subset is Borel. More
generally, every map from $\mathbb{Z}^2$ to any standard Borel space is Borel
measurable.
\end{lemma}

\begin{proof}
For each $p\in\mathbb{Z}^2$ the singleton $\{p\}$ is measurable (discrete
$\sigma$-algebra), and a countable union of measurable sets is measurable, so
every $A=\bigcup_{p\in A}\{p\}$ is Borel. For measurability of a map
$f\colon\mathbb{Z}^2\to Y$ with $Y$ standard Borel, it suffices that
$f^{-1}(B)$ be Borel for every Borel $B\subseteq Y$; since every subset of
$\mathbb{Z}^2$ is Borel, this holds trivially.
\end{proof}

\begin{theorem}[Borel chromatic number]
\label{thm:borel}
$\chi_B(G)=2$. In particular $\chi_B(G)\neq 5$.
\end{theorem}

\begin{proof}
\emph{Upper bound.} By Lemma~\ref{lem:all-borel} the parity coloring
$c_0\colon\mathbb{Z}^2\to\{0,1\}$ of Theorem~\ref{thm:bipartite} is Borel
measurable (its codomain $\{0,1\}$ is a standard Borel space), and it is proper.
Hence $\chi_B(G)\le 2$.

\emph{Lower bound.} As in Corollary~\ref{cor:chi}, $G$ has an edge, and any
proper coloring must assign distinct colors to the two endpoints of an edge, so
no $1$-coloring is proper. Therefore $\chi_B(G)\ge 2$.

Combining, $\chi_B(G)=2$.
\end{proof}

\begin{remark}[No extra color is forced]
Since $\chi(G)=\chi_B(G)=2$, the Borel constraint does not force an additional
color in this example. The phenomenon the conjecture invokes is real in other
settings---for instance, Borel graphs of bounded degree can have finite Borel
chromatic number strictly larger than their classical one---but it does not
occur for the square grid, whose bipartition is already Borel because the space
is countable and discrete.
\end{remark}

\section{Conclusion}

Both values asserted by conjecture 00000001367 are false:
\[
\chi(G)=2\quad(\text{not }4),\qquad \chi_B(G)=2\quad(\text{not }5).
\]
The Cayley graph of $\mathbb{Z}^2$ with generators $\{\pm e_1,\pm e_2\}$ is the
square grid, which is bipartite via the parity of $i+j$; and because
$\mathbb{Z}^2$ is countable and discrete, every subset---in particular the
parity classes---is Borel, so the very same $2$-coloring witnesses
$\chi_B\le 2$. Hence the conjecture is false in both of its numerical
assertions, and its explanation that ``the Borel constraint forces one extra
color'' does not apply here. The combinatorial part of this refutation is
mechanically verified in core Lean~4 (see \texttt{lean4/}); the Borel
measurability step is the standard argument of Lemma~\ref{lem:all-borel} and is
not formalised in Lean.

\end{document}
